% 分野別類題: リカッチ方程式（Riccati equation）解答例
% コンパイル: TEXINPUTS="../../../00_common:$TEXINPUTS" latexmk -xelatex answers.tex
\documentclass[a4paper,11pt]{article}
\usepackage{preamble}
\usepackage{shoakubox}

\begin{document}

\kamokumei{類題演習 解答例 --- Riccati Differential Equations}

\daimon{【解法】リカッチ方程式の解き方と導出}

\noindent
リカッチ方程式とは、次の形の1階微分方程式である。
\[
\frac{dy}{dx} = q_{0}(x) + q_{1}(x)\,y + q_{2}(x)\,y^{2}
\]
一般には求積法で解けないが、特解 $y_{1}(x)$ が1つ分かれば一般解が求められる。$y = y_{1} + \dfrac{1}{u}$ とおいて代入すると
\[
y_{1}' - \frac{u'}{u^{2}} = q_{0} + q_{1}\left(y_{1} + \frac{1}{u}\right) + q_{2}\left(y_{1} + \frac{1}{u}\right)^{2}
\]
$y_{1}$ が特解であること（$y_{1}' = q_{0} + q_{1}y_{1} + q_{2}y_{1}^{2}$）を用いて整理すると
\[
-\frac{u'}{u^{2}} = \frac{q_{1} + 2q_{2}y_{1}}{u} + \frac{q_{2}}{u^{2}}
\quad\Longleftrightarrow\quad
\frac{du}{dx} + \bigl(q_{1}(x) + 2q_{2}(x)\,y_{1}(x)\bigr)\,u = -q_{2}(x)
\]
という $u$ についての1階線形微分方程式に帰着する。これを積分因子法で解き、$y = y_{1} + \dfrac{1}{u}$ に戻せば一般解が得られる。

\clearpage
\begin{mondaiboxS}{問1}
$y_{1} = x$ が $\dfrac{dy}{dx} = 1 + x^{2} - 2xy + y^{2}$ の特解であることを確かめ、一般解を求めよ。
\end{mondaiboxS}
\kaitou
$y_{1} = x$ を代入すると、左辺 $= 1$、右辺 $= 1 + x^{2} - 2x^{2} + x^{2} = 1$ となり一致するので、$y_{1} = x$ は特解である。

$y = x + \dfrac{1}{u}$ とおくと $\dfrac{dy}{dx} = 1 - \dfrac{u'}{u^{2}}$。方程式に代入して
\[
1 - \frac{u'}{u^{2}} = 1 + x^{2} - 2x\left(x + \frac{1}{u}\right) + \left(x + \frac{1}{u}\right)^{2}
= 1 + \frac{1}{u^{2}}
\]
（$x^{2} - 2x^{2} - \dfrac{2x}{u} + x^{2} + \dfrac{2x}{u} = 0$ に注意。）よって
\[
-\frac{u'}{u^{2}} = \frac{1}{u^{2}}
\quad\Longleftrightarrow\quad
\frac{du}{dx} = -1
\quad\Longrightarrow\quad
u = C - x
\]
したがって一般解は
\[
\boxed{\; y = x + \frac{1}{C - x} \;}
\qquad (C \text{ は任意定数})
\]

\clearpage
\begin{mondaiboxS}{問2}
$y_{1} = \dfrac{1}{x}$ が $\dfrac{dy}{dx} = -\dfrac{1}{x^{2}} - \dfrac{y}{x} + y^{2}$ の特解であることを確かめ、一般解を求めよ。
\end{mondaiboxS}
\kaitou
$y_{1} = \dfrac{1}{x}$ を代入すると、左辺 $= -\dfrac{1}{x^{2}}$、右辺 $= -\dfrac{1}{x^{2}} - \dfrac{1}{x^{2}} + \dfrac{1}{x^{2}} = -\dfrac{1}{x^{2}}$ となり一致するので特解である。

$q_{1}(x) = -\dfrac{1}{x}$, $q_{2}(x) = 1$ であるから、$y = y_{1} + \dfrac{1}{u}$ の置換により $u$ は
\[
\frac{du}{dx} + \left(q_{1} + 2q_{2}y_{1}\right)u = -q_{2}
\quad\Longleftrightarrow\quad
\frac{du}{dx} + \frac{u}{x} = -1
\]
を満たす。積分因子は $\mu = x$ であるから
\[
\frac{d}{dx}(xu) = -x
\quad\Longrightarrow\quad
xu = -\frac{x^{2}}{2} + \frac{C}{2}
\quad\Longrightarrow\quad
u = \frac{C - x^{2}}{2x}
\]
したがって一般解は
\[
\boxed{\; y = \frac{1}{x} + \frac{2x}{C - x^{2}} \;}
\qquad (C \text{ は任意定数})
\]

\clearpage
\begin{mondaiboxS}{問3}
$y_{1} = e^{x}$ が $\dfrac{dy}{dx} = e^{-x}y^{2} + y - e^{x}$ の特解であることを用いて、一般解を求めよ。
\end{mondaiboxS}
\kaitou
（確認: $y_{1} = e^{x}$ のとき右辺 $= e^{-x}e^{2x} + e^{x} - e^{x} = e^{x} = y_{1}'$。）

$q_{1}(x) = 1$, $q_{2}(x) = e^{-x}$ であるから、$y = e^{x} + \dfrac{1}{u}$ の置換により
\[
\frac{du}{dx} + \left(1 + 2e^{-x}\cdot e^{x}\right)u = -e^{-x}
\quad\Longleftrightarrow\quad
\frac{du}{dx} + 3u = -e^{-x}
\]
積分因子は $\mu = e^{3x}$ であるから
\[
\frac{d}{dx}\left(u\,e^{3x}\right) = -e^{2x}
\quad\Longrightarrow\quad
u\,e^{3x} = -\frac{1}{2}e^{2x} + \frac{C}{2}
\quad\Longrightarrow\quad
u = \frac{C - e^{2x}}{2e^{3x}}
\]
したがって一般解は
\[
\boxed{\; y = e^{x} + \frac{2e^{3x}}{C - e^{2x}} \;}
\qquad (C \text{ は任意定数})
\]

\end{document}
